\subsection{An incoming loss and two matched lower roots}

An incoming lower loss can have unbounded transit time. To compare
its central derivative with that of a regular outgoing arc, we use
negative divergence on the outer arc and contraction through a fixed
Gaussian core. This argument also allows an incoming endpoint at
infinity. When both lower roots are lost, neither derivative provides
a regular reference; we instead compare their logarithmic rates in
the two moving deficits.

\begin{lemma}[Incoming contraction]\label{rw:finite:incoming-contraction}
Suppose the outgoing lower data are compact and regular, with
$A_+(p)\le A_{\rm out}$. Choose a fixed incoming crosscut $R>0$
with $F_->f_*>0$ on $[0,R]$ and entries $q>R+2\zeta$, where
$\zeta>0$ is fixed. On the actual common positive-central passing
set, for either sign of $\ell$, the incoming map $\mathcal I(q)$
satisfies
\begin{equation}\label{rw:finite:incoming-bound}
 0<-\mathcal I'(q)\le\frac C\epsilon e^{-A_-(R)/\epsilon}
        \left(\frac{1+R^2}{1+q^2}\right)^{1/(24\epsilon)}.
\end{equation}
If $A_-(R)>A_{\rm out}+3$ and the upper derivative is bounded,
then $K<0$ on that set. The closing function is strictly monotone
on each of its intervals. Every closing zero there is simple and
has return multiplier less than one.
\end{lemma}
\begin{proof}
In the incoming coordinates $W=-X$, the forward divergence on
$W\ge R$, $0<z<1$ is
\[
 -(1+2\epsilon\bar B)W
 +\epsilon\{\bar a+\bar c(2z-1)\}
                      -2\rho\ell(1-z)\le-W/2.
\]
The normal component satisfies
$z'\le-(R/2)z+\rho\max(\ell,0)$.
On the fixed strip $R\le W\le R+\zeta$ the horizontal speed
has magnitude at most $M\epsilon$. Any passing arc therefore
reaches $W=R$ with
\begin{equation}\label{rw:finite:height-at-crosscut}
 z_R\le e^{-R\zeta/(2M\epsilon)}
                          +2\rho\max(\ell,0)/R.
\end{equation}
An arc that has already reached $z=0$ is outside the stated domain.

The fixed core from this crosscut to the center has denominator
\[
 D_-=F_-+\epsilon\bar m-(1+2\epsilon\bar m)z
                                      +\epsilon\bar m z^2>f_*/2
\]
on a small height strip, and its graph equation is
\[
 z_W=\frac{Wz+\epsilon\bar c z(1-z)-\rho\ell(1-z)^2}
                  {\epsilon D_-}=:f(W,z,\ell).
\]
Let its signed output be $\mathcal C_R(z_R,\ell)$.
The Gaussian estimate gives
\[
 0\le\mathcal C_R(z_R,0)\le Cz_Re^{-A_-(R)/\epsilon}.
\]
For this bound, integrate
$\partial_W\log z=W/(\epsilon F_-)+O(1+(W/\epsilon+1)z)$;
the nonlinear integral is bounded by $Cz_R$.
For $\ell<0$ with positive central output, the solutions for
$\theta\ell$, $0\le\theta\le1$, lie between this positive
solution and the unforced one by scalar comparison. At fixed $z_R$,
variation of constants gives
\[
 \partial_\ell\mathcal C_R(z_R,\theta\ell)
 =\int_0^R e^{-\int_0^u f_z(v,z(v),\theta\ell)\,\dd v}
       \frac{(1-z(u))^2}{\rho D_-(u,z(u))}\,\dd u\ge c>0.
\]
To obtain the positive lower bound, restrict the integral to
$0\le u\le\rho$. Its prefactor is at least $c/\rho$, and
$|f_z|\le C(W/\epsilon+1+|\ell|/\rho)$ bounds the exponent
there. Positivity of the final output consequently implies
$|\ell|\le Cz_Re^{-A_-(R)/\epsilon}$.
For $\ell\ge0$, positivity of the outgoing regular map instead
gives $\ell\le CZ_+(p)$. Thus $\ell$ is exponentially small
on either part of the common passing set. Equation
\eqref{rw:finite:height-at-crosscut} now makes $z_R$ exponentially
small as well.

Set $\Phi(W)=\log(1+W^2)/(12\epsilon)$.
Since $W'=-\epsilon D_-$ and
$D_-\le2+W+W^2\le3(1+W^2)$, its time derivative is at least
$-W/2$, hence at least the divergence. The determinant formula
from $(q,1)$ to $W=R$ therefore bounds the transverse derivative
by
\[
 \frac{Cq}{\epsilon}
      \left(\frac{1+R^2}{1+q^2}\right)^{1/(12\epsilon)}.
\]
Half this power absorbs $q$ uniformly even for unbounded entries.
On the core, its normal variational exponent is
$A_-(R)/\epsilon+O(1)$, because
\[
 f_z=\frac{W}{\epsilon(F_-+\epsilon\bar m)}
       +\frac{\bar c}{F_-+\epsilon\bar m}
       +O((W/\epsilon+1)|z|+|\ell|/\rho),
\]
and the last two quantities have the bounds just proved.
Composition gives \eqref{rw:finite:incoming-bound}.
The outgoing relative amplitude has
$-L_+'(p)\ge(c/\epsilon)e^{-A_{\rm out}/\epsilon}$ on the
same compact regular tube, including the small forcing just bounded.
Thus the ratio of incoming and outgoing central derivatives is
less than $1/2$ when $A_-(R)>A_{\rm out}+3$ and the upper
derivative is bounded. This proves $K<0$ and the remaining claims.
The estimate uses divergence integrated along the actual passing arc
and remains valid for arbitrarily long residence near the incoming root.
\end{proof}

At a first finite simple or double incoming root, $A_-(R)\to\infty$
as $R$ approaches the root from below. If the outgoing arc and upper
return remain regular, choose $R$ so that the preceding lemma applies.
The resulting terminal part is monotone even though its incoming
transit time diverges.

For simultaneous finite lower losses,
Lemma~\ref{rw:finite:lower-cones} gives $\bar B=-\beta<0$ and
simple roots $p_l,q_l$ with $p_lq_l=1/\beta$. The next lemma
determines when a regular upper return joins these roots and supplies
the transverse parameter needed to measure their splitting.

\begin{lemma}[The regular two-root match]
\label{rw:finite:two-root-geometry}
For fixed $\beta>0$, $\bar a$ and $\bar m$, the equation
$\mathcal U_0(p_l;\bar m,\bar c)=q_l$ has a unique regular
solution in $\bar c$ exactly when $\bar m<\beta$.
At that solution the derivative with respect to $\bar c$ is
strictly positive. If $\bar a\ne0$, then
\begin{equation}\label{rw:finite:two-root-signs}
 \bar a\bar c<0,\qquad
 \operatorname{sign}(\bar a+\bar c)=\operatorname{sign}\bar c,
 \qquad |\bar c|>|\bar a|.
\end{equation}
For $\bar a=0$ the solution is $\bar c=0$.
\end{lemma}
\begin{proof}
On a returning upper arc $\bar mY-1<0$. Indeed the first
maximum of $Y$ must have $Y''=\bar mY-1<0$; the region below
that nullcline contains the complete arc. Its graph equation is
$Y_X=(X+\bar cY)/(\bar mY-1)$. Differentiation in $\bar c$,
with integration from $X=p_l$ in the decreasing direction, has
positive forcing after reversing the integration direction, because
$Y>0$ and $\bar mY-1<0$. It gives $Y_{\bar c}>0$ at the
exit and hence $\partial_{\bar c}\mathcal U_0>0$.
For $\bar m>0$, the formulas for $S,I$ show that the range at
fixed $p_l$ is $(0,1/(\bar mp_l))$: as the input approaches its
saddle threshold, the exit approaches the reflected threshold and
their product is $1/\bar m$. For $\bar m\le0$ that range is
$(0,\infty)$, by the same exponential, repeated and trigonometric
formulas, or their limits as $\bar c$ varies. Thus $q_l=1/(\beta p_l)$
belongs to the range exactly when $\bar m<\beta$.

To determine the coefficient signs at the match, put
$M=\beta-\bar m>0$ and $\tau=\bar a+\bar c$.
The endpoint identities from the solution $pS-I$ give
\[
 p_l+q_l=S(T)p_l(Mp_l-\tau)
       =e^{-\bar cT}S(T)q_l(Mq_l+\tau),
\]
and hence
\begin{equation}\label{rw:finite:matching-tau}
 \tau=M\frac{e^{\bar cT}p_l^2-q_l^2}
                      {e^{\bar cT}p_l+q_l}.
\end{equation}
The mechanical equation $Y''-\bar cY'-\bar mY=-1$ has energy
$\frac12(Y')^2+Y-\frac12\bar mY^2$ with derivative
$\bar c(Y')^2$. Thus $q_l-p_l$ has the sign of $\bar c$.
For the sign of \eqref{rw:finite:matching-tau}, write
$S(t)=e^{\bar ct/2}S_0(t)$. The positive fundamental solution
$S_0$ is hyperbolic, linear or trigonometric. For $T/2<u<T$
one has $S_0(u)>S_0(T-u)$: in the trigonometric case use
$T<\pi/\sqrt{-\bar m-\bar c^2/4}$ and the sine difference
formula. Pairing $u$ and $T-u$ then shows that
\[
 \int_0^T e^{\bar cu/2}S_0(u)\,\dd u
       -e^{\bar cT/2}\int_0^T e^{-\bar cu/2}S_0(u)\,\dd u
\]
has sign $\bar c$.
Since $q_l/p_l=e^{\bar cT}I_{-\bar c}(T)/I_{\bar c}(T)$,
this is precisely the sign of $e^{\bar cT}p_l^2-q_l^2$.
It proves the sign of $\tau$. Finally $p_l-q_l=\bar a/\beta$
gives \eqref{rw:finite:two-root-signs} and its zero case.
\end{proof}

Use the actual strong intersections of the two lower roots, their
positive regular power inputs, and the regular upper map. Common
finite-jet convergence preserves the preceding transverse derivative,
so their splitting is a coefficient coordinate $\delta$. The connector
then has the form $g(t)=(t-\delta)k(t)$, with $k,g'>c>0$ and
bounded selected jets on larger signed intervals. The deficit comparison
interval is $t>0$, $g(t)>0$; an actual return also requires positive
complete central outputs. As in Section~5, extend each absolute map by
zero on its own nonpassing side.

\begin{proposition}[Two-root local count]
\label{rw:finite:two-root-count}
Near every regular two-root match there are common coefficient
and section neighborhoods and a scale cutoff such that the closing
function has at most three isolated zeros, including multiplicity,
on its positive central passing interval. The corresponding
logarithmic ratio has at most two zeros and its derivative at most
one. These assertions include either sign of splitting and the
reversible match.
\end{proposition}
\begin{proof}
The exact function is
\[
 F=D_+(t)-D_-(g(t))-\rho\nu H,\qquad
 D_\pm(x)=K_\pm x^{h_\pm/\epsilon}e^{j_\pm(x)},
\]
where $h_+=p_l/\Delta+O(\epsilon)$,
$h_-=q_l/\Delta+O(\epsilon)$ and $\Delta=\sqrt{\bar a^2+4\beta}$.
All reference terms are input independent. Put
\[
 R_\pm=j_\pm'+\epsilon
       \frac{j_\pm'+xj_\pm''}{h_\pm+\epsilon xj_\pm'}.
\]
The denominator has a positive margin. Thus
$D_\pm''/D_\pm'=(h_\pm/\epsilon-1)/x+R_\pm$
and the exact cleared numerator is
\begin{equation}\label{rw:finite:two-root-numerator}
 M=tgK'=(h_--\epsilon)tg'-(h_+-\epsilon)g
        +\epsilon tg\{g'R_-(g)-R_+(t)+g''/g'\}.
\end{equation}
The multiplier $tg$ is positive on the whole passing interval.
Away from $\bar a=0$,
$h_--h_+=-\bar a/(\beta\Delta)+O(\epsilon)$ has a nonzero
margin. Differentiating \eqref{rw:finite:two-root-numerator} gives
\[
 M'=(h_--h_+)g'+(h_--\epsilon)tg''+O(\epsilon(t+g)).
\]
First choosing the section radius small, then the coefficient
neighborhood and scale small, makes $M'$ have a fixed nonzero
sign. Thus $M$ has at most one simple zero there.

Near the reversible match use the independent transverse variables
$(\bar a,\delta,\nu)$. The inverse coefficient chart follows
from the strict splitting derivative just proved; its finite jets
are bounded on common boxes. At their zero, the reflected power
coordinates may be chosen equal, giving $D_+=D_-$ and $g=t$.
The relative coefficients extend by the common linear rule of
Section~5, so the required coefficient segments are legitimate.
At zero scale the columns of the numerator are
\[
 M_{\bar a}=-\frac{t}{2\beta^{3/2}}+O(t^2),\qquad
 M_\delta=\frac1{2\beta}+O(t^2).
\]
Indeed $p_l=q_l=1/\sqrt\beta$, $\Delta=2\sqrt\beta$ and
$(h_--h_+)_{\bar a}=-1/(2\beta^{3/2})$; differentiating
$g=(t-\delta)k$ with the common reflected unit $k(0)=1$
gives the constant column. The same formulas hold after a common
positive normalization if a different reflected unit was used.
Consequently the second column has a positive margin and the
derivative of their quotient is strictly negative on a short
interval including zero. One coefficient Hadamard integral now
gives
\[
 M=\bar a A_\epsilon+\delta B_\epsilon+\nu C_\epsilon,
 \qquad B_\epsilon>c,\quad (A_\epsilon/B_\epsilon)'<-c.
\]
The finite-order preparations reserve two ordinary input derivatives
for each logarithmic rate and one coefficient derivative for this
division, which gives the displayed derivative bounds.

On the common small lower input intervals, the two large powers
and their required absolute coefficient derivatives are bounded
by $C\epsilon^{-N}e^{-b/\epsilon}$ with a fixed $b>0$.
Their exact reversible equality and the signed extensions give,
at any closing zero, $|\nu|=o(1)(|\bar a|+|\delta|)$ by
Lemma~\ref{rw:centers:relative-forcing}. If the latter two
parameters vanish, the fibre is a center and has no isolated zero.
Otherwise the strict quotient margin and relative $C^1$ forcing bound
give at most one simple zero of $M$ by
Lemma~\ref{rw:atlas:two-columns}.
Two Rolle steps give the claimed bounds for $K$ and $F$.
Both full central outputs increase with $t$, so their positivity
truncates the deficit interval to one interval. Excluded endpoints
contribute no formal closing zero.
\end{proof}

With a finite transverse upper limit, the lower possibilities are now
exhausted: both arcs are regular, the incoming arc alone loses
positivity, or both arcs lose positivity.
Lemma~\ref{rw:finite:action-order} excludes an outgoing loss alone.
The remaining finite upper limit is a saddle and is treated next.
